\chapter{Arithmetic coefficients and local conditions}
\label{v4-arith:coefficients:chapter}

A global cohomology class can satisfy a local condition because its
restriction differs from an allowed local class by a coboundary. The
primitive of that coboundary belongs to the relative cocycle. Retaining
it gives an explicit complex, and compatible local multiplications give
the two products of Chapter~\ref{v4-arith:gluing:chapter}. Forgetting it
before taking products loses part of the arithmetic input.

Coefficients must also be retained before taking limits. Every finite
\(\ell\)-power torsion group vanishes after tensoring with
\(\mathbb Q_\ell\), whereas their inverse limit can have positive
\(\mathbb Z_\ell\)-rank. We therefore construct the finite coefficient
complexes first, specify their transition maps, and only then complete.
Deformation functors impose a further requirement: all local and global
objects must be defined on the same coefficient category. Their tangent
complexes give local conditions on cochains. Their geometric
representability is a separate question.

\section{The arithmetic curve and its coefficients}
\label{v4-arith:coefficients:sec:places}

Fix a number field \(K\), its ring of integers \(\mathcal O_K\), and a
rational prime \(\ell\). A finite place \(v\) is a nonzero prime ideal
of \(\mathcal O_K\). Its completion \(K_v\) has valuation ring
\(\mathcal O_v\), maximal ideal \(\mathfrak m_v\), and finite residue
field \(k_v\) of characteristic \(p_v\). The arithmetic residue prime
\(p_v\) and the coefficient prime \(\ell\) serve different purposes.
They may coincide, but no such equality is imposed here.

The scheme \(\operatorname{Spec}\mathcal O_K\) has dimension one.
Its closed points are the finite places. Its generic point has field
\(K\). Archimedean completions are additional data, not additional
closed points of this scheme.

The adele ring \(\mathbb A_K\) is the
restricted product of the completions, with respect to \(\mathcal O_v\)
at finite places. An adele is a tuple of local elements. Identifying
such a tuple with a coordinate for chiral insertions would require a
geometric construction and its transition maps.

Let \(S\) be a finite set of places containing the archimedean places
and the places over \(\ell\). Write \(S_f\) for its finite part.
Let \(G_{K,S}\) be the Galois group of the maximal extension of \(K\)
unramified outside \(S\), and let \(G_{K_v}\) be the absolute Galois
group of \(K_v\).

Choose embeddings of algebraic closures defining
continuous homomorphisms
\[
 j_v:G_{K_v}\longrightarrow G_{K,S}\qquad(v\in S_f).
\]
These homomorphisms need not be injective. Strict restriction maps use
the chosen \(j_v\). Conjugate choices require cochain homotopies when
the resulting complexes are compared.

Put \(R_r=\mathbb Z/\ell^r\mathbb Z\), for \(r\geq1\).
Complexes are cohomological, with differential of degree one and
\(M[1]^n=M^{n+1}\). Tensor products are over the displayed coefficient
ring and use the Koszul sign rule. An ordinary tensor product denotes
the algebraic tensor product. A derived or completed tensor product
will be specified explicitly.

\section{Continuous cochains with their multiplication}
\label{v4-arith:coefficients:sec:cochains}

Let \(G\) be profinite, \(R\) a commutative ring with trivial
\(G\)-action, and \(M\) a discrete continuous \(R[G]\)-module.
The normalized continuous cochains \(C^n(G,M)\) are the continuous
functions \(G^n\to M\) vanishing when any argument is the identity.
Set \(C^0(G,M)=M\) and \(C^n(G,M)=0\) for \(n<0\).

Their differential is
\begin{align}
 (dc)(g_1,\ldots,g_{n+1})={}&g_1c(g_2,\ldots,g_{n+1})\nonumber\\
 &+\sum_{i=1}^n(-1)^i
 c(g_1,\ldots,g_ig_{i+1},\ldots,g_{n+1})\nonumber\\
 &+(-1)^{n+1}c(g_1,\ldots,g_n).
 \label{v4-arith:coefficients:eq:differential}
\end{align}
In degree zero this says \((dm)(g)=gm-m\).
Continuous group cohomology is the cohomology of this complex.\footnote{
J. S. Milne, \emph{Class Field Theory}, version 4.03 (2020),
Chapter II, \S4, pp.~86--89, gives the continuous-cochain and
finite-quotient conventions.}

Suppose now that \(M\) is an associative unital \(R\)-algebra and
that \(G\) acts by \(R\)-algebra automorphisms. For cochains \(c,c'\)
of degrees \(n,m\), define the ordered cup product by
\begin{equation}\label{v4-arith:coefficients:eq:cup}
 (c\smile c')(g_1,\ldots,g_{n+m})
 =c(g_1,\ldots,g_n)
   (g_1\cdots g_n)c'(g_{n+1},\ldots,g_{n+m}).
\end{equation}
No commutativity of \(M\) is needed. When the action is trivial, both
action factors in \eqref{v4-arith:coefficients:eq:differential}
and \eqref{v4-arith:coefficients:eq:cup} disappear.

\begin{proposition}\label{v4-arith:coefficients:prop:cochains}
The differential squares to zero. Under the algebra hypothesis,
\(C^\bullet(G,M)\) is a unital associative differential graded
\(R\)-algebra. Restriction along continuous group homomorphisms,
with the induced coefficient action, preserves its differential,
product, and unit.
\end{proposition}
\begin{proof}
Write \(d=\sum_i(-1)^i d_i\) for the faces in
\eqref{v4-arith:coefficients:eq:differential}. For \(i<j\), the
identity \(d_jd_i=d_id_{j-1}\) pairs two terms of \(d^2\) with
opposite signs. The identity for the first two faces uses
\(g(hm)=(gh)m\). The adjacent inner faces use associativity in
\(G\). Every term is paired, so \(d^2=0\). If an argument is the
identity, the two neighboring faces cancel and all remaining terms
vanish by normalization.

In the differential of a cup product, split the faces at the boundary
of its two argument blocks. The faces in the first block give
\((dc)\smile c'\), and those in the second give
\((-1)^n c\smile dc'\). The extra faces at the boundary cancel.
The action on the second factor in
\eqref{v4-arith:coefficients:eq:cup} makes this cancellation valid
also for a nontrivial coefficient action. This proves Leibniz's rule.

For three cochains the two bracketings evaluate the same three blocks.
The second value is acted on by the first block, and the third by
the first two blocks. The action is by algebra
automorphisms, so these expressions agree. The element \(1_M\) in
degree zero is a closed two-sided unit. Substitution of a group
homomorphism in each argument respects every displayed formula.
\end{proof}

\begin{proposition}\label{v4-arith:coefficients:prop:reduction}
For trivial coefficients \(R_r\), reduction gives degreewise
surjections and isomorphisms
\[
 C^n(G,R_{r+1})\otimes_{R_{r+1}}R_r
 \cong C^n(G,R_r).
\]
These maps preserve the differential, cup product, and restrictions.
Each module \(C^n(G,R_r)\) is flat over \(R_r\).
\end{proposition}
\begin{proof}
A continuous function from a compact space to a discrete space has
finite image and is constant on a finite clopen partition. Lift each
value from \(R_r\) to \(R_{r+1}\), choosing the zero lift for the zero
value. This preserves normalization.

A function reducing to zero has
values in \(\ell^rR_{r+1}\). Choosing preimages of its finitely many
values shows that the kernel is
\(\ell^rC^n(G,R_{r+1})\). This proves the tensor assertion.
All operations use only coefficient addition and multiplication.

Fix a finite clopen partition of \(G^n\). Normalized functions
constant on its cells form a finite free \(R_r\)-module. A basis
consists of the characteristic functions of cells disjoint from all
tuples with an identity argument. Refinement gives a
filtered system whose union is all normalized continuous cochains.
Filtered colimits of flat modules are flat.
\end{proof}

The finite models are compatible with the profinite topology. For
trivial coefficients, every continuous cochain factors through
\((G/U)^n\) for some open normal subgroup \(U\subset G\).
Indeed, its finite clopen fibers admit a finite cover by products of
cosets of open subgroups. A common open normal subgroup refining those
subgroups makes the function constant on the required product cosets.
Normalization then holds on \(G/U\), because a coset of the identity
has the same value as an actual identity argument.

Consequently
\[
 C^\bullet(G,R_r)=\varinjlim_U C^\bullet(G/U,R_r),
\]
with inflation as transition maps. The differential and cup product
are compatible with this expression.

This flatness assertion concerns modules in individual degrees.
It does not prove that an unbounded complex is \(K\)-flat, meaning
that tensoring it with any acyclic complex remains acyclic.
The resolution hypotheses of
Proposition~\ref{v4-prop:bar-derived-tensor} are still needed when
the coefficient bar is used to compute a derived tensor product.

\section{Local conditions and retained primitives}
\label{v4-arith:coefficients:sec:local}

Fix a discrete continuous \(G_{K,S}\)-module \(M\) over \(R\).
A cochain local condition at \(v\) consists of a complex \(L_v\)
and a chain map
\[
 g_v:L_v\longrightarrow C^\bullet(G_{K_v},M).
\]

Set \(A=C^\bullet(G_{K,S},M)\),
\(B=\prod_{v\in S_f}C^\bullet(G_{K_v},M)\), and
\(L=\prod_{v\in S_f}L_v\). Let \(f=(j_v^*)_v\) and
\(g=(g_v)_v\). The products are finite, so they are also direct sums
of complexes. Define the Selmer complex for these cochain data by
\begin{equation}\label{v4-arith:coefficients:eq:selmer-complex}
 \begin{split}
 C^n_{\mathcal L}&=A^n\oplus L^n\oplus B^{n-1},\\
 D(a,c,y)&=(da,dc,f(a)-g(c)-dy).
 \end{split}
\end{equation}
The symbol \(\mathcal L\) records the chosen complexes and maps.
A subspace of local cohomology alone does not specify these data.

\begin{proposition}\label{v4-arith:coefficients:prop:selmer-complex}
The complex \(C^\bullet_{\mathcal L}\) is the homotopy fiber of
\(f-g:A\oplus L\to B\). In particular it has the exact sequence
\[
 \cdots\longrightarrow H^{n-1}(B)
 \longrightarrow H^n(C^\bullet_{\mathcal L})
 \longrightarrow H^n(A)\oplus H^n(L)
 \xrightarrow{f_*-g_*}H^n(B)\longrightarrow\cdots.
\]
A cocycle is a pair of global and allowed local cocycles together
with a specified primitive of their difference after restriction.
\end{proposition}
\begin{proof}
The third component of \(D^2\) is
\(f(da)-g(dc)-d(f(a)-g(c)-dy)=0\).
Projection onto \(A\oplus L\) has kernel \(B[-1]\), with
differential \(-d_B\). The resulting short exact sequence is the
mapping-fiber sequence for \(f-g\), and its connecting map is
\(f_*-g_*\). The cocycle equations are exactly those stated.
\end{proof}

Three choices exhibit the range of this definition. The zero
condition \(L_v=0\) asks for a primitive of the local restriction.
The full condition \(L_v=C^\bullet(G_{K_v},M)\), with \(g_v\) the
identity, imposes no restriction. If every condition is full, the map
\(a\mapsto(a,f(a),0)\) is a quasi-isomorphism
\(A\to C^\bullet_{\mathcal L}\): its complementary fiber is the
contractible fiber of \(1_B\).

For a finite place, write \(I_v\) for inertia in \(G_{K_v}\).
The unramified cochain condition is
\begin{equation}\label{v4-arith:coefficients:eq:unramified}
 L_v=C^\bullet(G_{K_v}/I_v,M^{I_v})
 \xrightarrow{\mathrm{inflation}} C^\bullet(G_{K_v},M).
\end{equation}
The quotient acts on the inertia invariants, and inflation followed
by their inclusion defines the map. If \(M\) is a coefficient algebra
with Galois action by automorphisms, then \(M^{I_v}\) is a unital
subalgebra and inflation preserves the cup product.

General local conditions are complexes. When \(A,L,B\) are strict
dg algebras and \(f,g\) are algebra maps, there are additional products.
For \(x=(a,c,y)\) of degree \(n\), they are
\begin{align}
 \mu_L(x,x')&=(aa',cc',(-1)^nf(a)y'+yg(c')),
 \label{v4-arith:coefficients:eq:left-product}\\
 \mu_R(x,x')&=(aa',cc',yf(a')+(-1)^ng(c)y').
 \label{v4-arith:coefficients:eq:right-product}
\end{align}
Proposition~\ref{v4-arith:gluing:prop:products} proves their dg
identities. Chapter~\ref{v4-arith:gluing:chapter} constructs the
comparison in every tensor length, including the module and bimodule
maps. These conclusions apply to a local condition only after its
algebra structure and the multiplicativity of \(g_v\) have been given.

If the cospan is unital, \(u=(1_A,1_L,0)\) is the native unit of
both products. For the zero local condition, the displayed triple
provides no such common unit.

In either case the external unitization
\[
 R\oplus C^\bullet_{\mathcal L},\qquad
 (r,x)(r',x')=(rr',rx'+r'x+xx'),
\]
has augmentation \((r,x)\mapsto r\) and new unit \((1,0)\).
When \(u\) exists, \((0,u)\) is an idempotent in its augmentation
ideal. It is not the new unit. The coefficient bar uses this external
augmentation. Local vacua, scalar augmentations, and empty geometric
cuts retain their distinct definitions.

A comparison of cochain models may respect restriction only up to
homotopy. For the zero local condition, suppose algebra maps
\(\alpha:A\to A'\), \(\beta:B\to B'\) and a degree-minus-one map
\(H:A\to B'\) satisfy
\(f'\alpha-\beta f=dH+Hd\). Then
\[
 (a,y)\longmapsto(\alpha(a),\beta(y)+H(a))
\]
is a chain map of relative complexes. Its multiplicative defect is
computed in Proposition~\ref{v4-arith:relative-bar:prop:homotopy}.
A nonzero defect requires higher comparison maps before one obtains
a bar-coalgebra map. A chain homotopy alone is not that construction.

\section{Completion before rationalization}
\label{v4-arith:coefficients:sec:completion}

Give \(\mathbb Z_\ell=\varprojlim_rR_r\) its inverse-limit topology.
A map \(G^n\to\mathbb Z_\ell\) is continuous exactly when each of
its reductions is continuous. Hence normalized cochains satisfy
\[
 C^\bullet_{\mathrm{cts}}(G,\mathbb Z_\ell)
 =\varprojlim_r C^\bullet(G,R_r).
\]
The differential, cup product, and restrictions pass to this limit.
For a general family of local conditions \(L_{v,r}\), transition
maps and coefficient base-change properties must be supplied as part
of that family.

\begin{proposition}\label{v4-arith:coefficients:prop:milnor}
Let \(E_r\) be a countable tower of complexes with degreewise
surjective transition maps. There is an exact sequence
\[
 0\longrightarrow\varprojlim{}^1_r H^{n-1}(E_r)
 \longrightarrow H^n(\varprojlim_r E_r)
 \longrightarrow\varprojlim_r H^n(E_r)\longrightarrow0.
\]
Here \(\varprojlim{}^1\) is the cokernel of \(1-\mathrm{shift}\)
on the product of the indicated cohomology groups.
\end{proposition}
\begin{proof}
The map \(1-\mathrm{shift}:\prod_rE_r\to\prod_rE_r\) is
surjective in every degree. To solve \(x_r-t_r(x_{r+1})=y_r\),
choose \(x_1\) and lift \(x_r-y_r\) recursively along the
surjection \(t_r:E_{r+1}\to E_r\). Its kernel is the inverse limit.
Thus
\[
 0\longrightarrow\varprojlim_rE_r\longrightarrow\prod_rE_r
 \xrightarrow{1-\mathrm{shift}}\prod_rE_r\longrightarrow0
\]
is an exact sequence of complexes. Products of modules are exact,
so their cohomology is the product of the cohomology groups.
The resulting long exact sequence gives the assertion.
\end{proof}

The hypothesis holds for the coefficient cochains and their relative
complexes with zero or full local conditions. It must be checked for
other conditions. Even under this hypothesis, the possible
\(\varprojlim{}^1\) term prevents replacing the cohomology of the
limit by the limit of the cohomology groups without further argument.

For a compatible tower of relative dg algebras \(P_r\) over \(R_r\),
let \(\mathcal B_r\) be the coefficient bar of \(R_r\oplus P_r\).
It is the direct sum of its finite tensor words, with the signs of
\eqref{v4-eq:bar-coefficient-signs}. Its coefficient completion and
coproduct target are
\begin{equation}\label{v4-arith:coefficients:eq:completed-bar}
 \widehat{\mathcal B}=\varprojlim_r\mathcal B_r,
 \qquad
 \widehat{\mathcal B}\,\widehat\otimes\,
 \widehat{\mathcal B}
 :=\varprojlim_r(\mathcal B_r\otimes_{R_r}\mathcal B_r).
\end{equation}
Every element has finite word support at each coefficient precision.
This is smaller than an unrestricted product over word lengths.
The differential, coproduct, and endpoint comparisons pass to
\eqref{v4-arith:coefficients:eq:completed-bar} by their compatible
finite-stage formulas, as proved in
Chapters~\ref{v4-ch:bar-cobar-topology}
and~\ref{v4-arith:gluing:chapter}.

Each finite-stage bar is conilpotent: sufficiently many reduced cuts
kill any finite sum of words. The limit need not be conilpotent.
Take a free rank-one module with basis \(a\), where \(|a|=1\).
The degree-zero
element \(\sum_{n\geq1}\ell^n(sa)^{\otimes n}\) has finite support
modulo every \(\ell^r\). Every iterated reduced coproduct has a
nonzero term. Its existence depends on coefficient completion, not on a
bound on geometric configurations.

\begin{example}\label{v4-arith:coefficients:ex:rationalization}
Each \(R_r\otimes_{\mathbb Z_\ell}\mathbb Q_\ell\) is zero, since
\(\ell^r\) annihilates its first factor and is invertible on the
second. But
\[
 (\varprojlim_rR_r)\otimes_{\mathbb Z_\ell}\mathbb Q_\ell
 =\mathbb Q_\ell.
\]
The constant degree-zero cochains already exhibit this difference.
There is no unital ring map \(R_r\to\mathbb Q_\ell\): the relation
\(\ell^r1=0\) would persist in the target. Thus a characteristic-zero
realization cannot be constructed by rationalizing these finite
coefficient algebras separately.
\end{example}

\section{Restriction, transfer, and tensor support}
\label{v4-arith:coefficients:sec:transfer}

Restriction and coefficient reduction are multiplicative in the models
above. Additive transfer has a different behavior. If \(H\subset G\)
has finite index \(n\), corestriction on degree-zero cohomology with
trivial \(\mathbb Z_\ell\)-coefficients is multiplication by \(n\).
For finite groups this is the coset sum. For profinite groups the
same formula uses the finite coset space.

\begin{proposition}\label{v4-arith:coefficients:prop:transfer}
If \(n>1\), corestriction with trivial \(\mathbb Z_\ell\)-coefficients
is not multiplicative. There is no \(A_\infty\) morphism of these
cochain algebras whose induced cohomology map is corestriction.
\end{proposition}
\begin{proof}
On the class of the degree-zero unit, multiplicativity would give
\(n=\operatorname{cor}(1\cdot1)
=\operatorname{cor}(1)\operatorname{cor}(1)=n^2\).
The ring \(\mathbb Z_\ell\) has characteristic zero, so this fails
for \(n>1\). The arity-two identity for an \(A_\infty\) morphism
makes its first component multiplicative on cohomology, yielding
the same contradiction. For index \(\ell\), reduction modulo
\(\ell^2\) already detects the failure: \(\ell\ne\ell^2=0\).
\end{proof}

Multiplicative norm operations, when constructed, have different
axioms and cannot be substituted for this additive transfer without
specifying their coefficient and algebra structures.

Changing the coefficient prime creates another separation. An abelian
group is \(\ell\)-primary torsion if every element is killed by some
power of \(\ell\). For different primes, the derived tensor of two
such groups vanishes.

\begin{proposition}\label{v4-arith:coefficients:prop:primary}
If \(M_\ell\) and \(M_p\) are respectively \(\ell\)-primary and
\(p\)-primary torsion groups, with \(\ell\ne p\), then
\(M_\ell\otimes^{\mathbf L}_{\mathbb Z}M_p=0\).
For equal primes the corresponding assertion is false.
\end{proposition}
\begin{proof}
The resolution
\(0\to\mathbb Z\xrightarrow{\ell^a}\mathbb Z
\to\mathbb Z/\ell^a\to0\), tensored with \(M_p\), is acyclic:
multiplication by \(\ell^a\) is invertible on every \(p\)-primary
group. Every finitely generated subgroup of \(M_\ell\) is a finite
sum of cyclic \(\ell\)-power groups. Filtered unions and their exact
tensor and Tor functors extend the calculation to \(M_\ell\).
There is no higher Tor over \(\mathbb Z\). For equal primes,
\[
 \mathbb F_\ell\otimes^{\mathbf L}_{\mathbb Z}\mathbb F_\ell
 \simeq\mathbb F_\ell\oplus\mathbb F_\ell[1],
\]
since the differential in the same two-term resolution becomes zero.
The shift puts its Tor group in cohomological degree \(-1\).
\end{proof}

This calculation concerns tensor products over \(\mathbb Z\).
It neither supplies comparison maps between local geometric categories
nor makes two factors with the same coefficient field disjoint.

\section{Deformation functors on one coefficient category}
\label{v4-arith:coefficients:sec:deformations}

Local conditions on representations supply cochain conditions by
linearization. To form their global pullback, first fix a finite field
\(k\) of characteristic \(\ell\), its Witt ring \(W(k)\), and a
continuous residual representation
\[
 \bar\rho:G_{K,S}\longrightarrow\mathrm{GL}_d(k).
\]
Use a single category \(\mathrm{Art}_{W(k)}\) of augmented simplicial
Artinian local \(W(k)\)-algebras with residue field \(k\).
Deformation functors take values in spaces and invert weak
equivalences of these coefficient algebras. A local theory over the
Witt ring of a different residue characteristic is not automatically
a functor on this category.

For a profinite group \(G\) with the specified residual representation,
write \(\mathrm{Def}_G(A)\) for its derived continuous deformation
space. Concretely, it is the homotopy fiber at \(\bar\rho\) of
\[
 \operatorname{Map}_{\mathrm{cts}}(BG,B\mathrm{GL}_d(A))
 \longrightarrow
 \operatorname{Map}_{\mathrm{cts}}(BG,B\mathrm{GL}_d(k)),
\]
using the continuous derived mapping-space model.

Here \(BG\) is the classifying space of \(G\), and the fiber retains
an identification of the reduction with \(\bar\rho\). For ordinary
Artinian \(A\), objects are lifts of \(\bar\rho\). Arrows are
conjugations by matrices reducing to the identity. These arrows retain
infinitesimal automorphisms.

Keeping a basis and forbidding those
conjugations gives the framed deformation problem instead.\footnote{
S. Galatius and A. Venkatesh, \emph{Derived Galois deformation rings},
arXiv:1608.07236, \S\S4--5 and 7, construct the derived continuous
deformation functors and their tangent complexes. Their \S9 imposes
local conditions by homotopy pullbacks.}

Let \(\mathrm{ad}\bar\rho=M_d(k)\), with action
\(g\cdot X=\bar\rho(g)X\bar\rho(g)^{-1}\).
The tangent complex of a deformation functor records its linear
response to square-zero coefficient extensions. With the cohomological
convention here, its degrees \(-1,0,1\) describe infinitesimal
automorphisms, deformations, and obstructions.

\begin{proposition}\label{v4-arith:coefficients:prop:tangent}
The tangent complex of \(\mathrm{Def}_G\) at \(\bar\rho\) is
\[
 T\mathrm{Def}_G\simeq C^\bullet(G,\mathrm{ad}\bar\rho)[1].
\]
In particular the unframed first-order deformation space is
\(H^1(G,\mathrm{ad}\bar\rho)\), its infinitesimal automorphisms are
\(H^0(G,\mathrm{ad}\bar\rho)\), and the framed first-order space is
\(Z^1(G,\mathrm{ad}\bar\rho)\).
\end{proposition}
\begin{proof}
Over \(k[\epsilon]/\epsilon^2\), write a lift as
\(\rho_\epsilon(g)=(1+\epsilon c(g))\bar\rho(g)\).
The representation equation becomes
\[
 c(gh)=c(g)+g\cdot c(h),
\]
the cocycle equation in
\eqref{v4-arith:coefficients:eq:differential}.
Conjugation by \(1+\epsilon X\) changes \(c\) by
\(X-g\cdot X=-dX\). Quotienting gives \(H^1\), while the
stabilizer is \(\ker d=H^0\). Before quotienting, the lifts form
\(Z^1\).

The same calculation applies to a simplicial square-zero module \(N\).
The kernel of
\(\mathrm{GL}_d(k\oplus N)\to\mathrm{GL}_d(k)\) is the additive
simplicial module \(M_d(N)\), since products of two elements of
\(N\) vanish. Conjugation induces the adjoint action.

The continuous cosimplicial mapping model has normalized cochains
\(C^\bullet(G,\mathrm{ad}\bar\rho\otimes_k N)\).
Passing to the classifying space shifts this linear complex by one.
This identifies the linear response on square-zero modules, hence
the stated tangent complex. Automorphisms remain in degree \(-1\).
\end{proof}

Fix a determinant lift \(\delta:G_{K,S}\to W(k)^\times\).
There are two distinct meanings of fixed determinant. One may require
the determinant character to have the value \(\delta\), while allowing
all strict \(\mathrm{GL}_d\)-conjugations. Alternatively, one may retain
an identification of the determinant representation with \(\delta\)
and require arrows to preserve that identification. The second problem
is the homotopy fiber of the determinant map of deformation functors.
Denote it by \(\mathrm{Def}_{G,\delta}^{\mathrm{or}}\).

\begin{proposition}\label{v4-arith:coefficients:prop:determinant}
For the determinant-identified problem,
\[
 T\mathrm{Def}_{G,\delta}^{\mathrm{or}}
 \simeq C^\bullet(G,\mathrm{ad}^0\bar\rho)[1],
 \qquad
 \mathrm{ad}^0\bar\rho=\ker(\operatorname{tr}:M_d(k)\to k).
\]
This holds also when \(\ell\) divides \(d\).
\end{proposition}
\begin{proof}
For a square-zero matrix, expansion of the determinant gives
\(\det(1+\epsilon X)=1+\epsilon\operatorname{tr}X\).
The tangent map of determinant is therefore the trace map
\[
 C^\bullet(G,\mathrm{ad}\bar\rho)[1]
 \longrightarrow C^\bullet(G,k)[1].
\]

Trace is equivariant and surjective. It is surjective on normalized
cochains in every degree: lift a scalar value \(a\) to
\(\operatorname{diag}(a,0,\ldots,0)\). This choice preserves
continuity and normalization, though it need not be equivariant as
a splitting. The kernel complex is
\(C^\bullet(G,\mathrm{ad}^0\bar\rho)\), so the homotopy fiber is
quasi-isomorphic to that kernel. No division by \(d\) occurs.
\end{proof}

For the determinant-value problem modulo all strict conjugations,
the first-order quotient is instead
\[
 Z^1(G,\mathrm{ad}^0\bar\rho)\big/d\bigl(M_d(k)\bigr),
\]
and its infinitesimal automorphisms are \(H^0(G,\mathrm{ad}\bar\rho)\).
For a trivial two-dimensional representation, these automorphisms
have dimension four. Preserving a determinant identification leaves
the trace-zero automorphisms, of dimension three, also in
characteristic two. The notation for the deformation problem must
therefore specify which arrows are permitted.

\section{The Selmer pullback and its tangent}
\label{v4-arith:coefficients:sec:selmer-functor}

Use the determinant-identified convention throughout this section.
Write \(\mathcal D=\mathrm{Def}_{G_{K,S},\delta}^{\mathrm{or}}\)
and \(\mathcal D_v=\mathrm{Def}_{G_{K_v},\delta|_v}^{\mathrm{or}}\).
Choose pointed local deformation functors \(\mathcal D_v^{\mathcal L}\)
and actual morphisms \(\mathcal D_v^{\mathcal L}\to\mathcal D_v\),
all on \(\mathrm{Art}_{W(k)}\). Define
\begin{equation}\label{v4-arith:coefficients:eq:selmer-functor}
 \mathcal D_{\mathcal L}
 =\mathcal D\times^h_{\prod_{v\in S_f}\mathcal D_v}
                    \prod_{v\in S_f}\mathcal D_v^{\mathcal L}.
\end{equation}
The symbol \(\times^h\) denotes a homotopy pullback of functors to
spaces. It imposes compatibility together with its homotopy, rather
than equality only on isomorphism classes.

\begin{theorem}\label{v4-arith:coefficients:thm:selmer-tangent}
Suppose each chosen local functor has tangent complex \(L_v[1]\),
and its morphism to \(\mathcal D_v\) induces a chain map
\(g_v:L_v\to C^\bullet(G_{K_v},\mathrm{ad}^0\bar\rho)\).
Then
\[
 T\mathcal D_{\mathcal L}
 \simeq C^\bullet_{\mathcal L}
          (G_{K,S},\mathrm{ad}^0\bar\rho)[1],
\]
where the complex on the right is
\eqref{v4-arith:coefficients:eq:selmer-complex} with these actual
maps. Geometric representability and Noetherianity are additional
conditions on this formal deformation functor.
\end{theorem}
\begin{proof}
Evaluate \eqref{v4-arith:coefficients:eq:selmer-functor} on every
square-zero coefficient extension and take the fiber at its base
point. Homotopy limits commute with homotopy limits, so this operation
preserves the displayed pullback. The associated tangent spectra,
and thus their complexes, give the fiber of
\[
 C^\bullet(G_{K,S},\mathrm{ad}^0\bar\rho)[1]
 \oplus\prod_vL_v[1]
 \longrightarrow
 \prod_vC^\bullet(G_{K_v},\mathrm{ad}^0\bar\rho)[1].
\]
The arrow is restriction minus the chosen local map. Its mapping
fiber is exactly \eqref{v4-arith:coefficients:eq:selmer-complex},
shifted by one.
\end{proof}

If each \(g_v\) is an inclusion of complexes, the quotient complex
\(C^\bullet(G_{K_v},\mathrm{ad}^0\bar\rho)/L_v\) gives an equivalent
description of the same fiber. Without such inclusions the quotient
notation has no supplied meaning. The mapping cone remains valid.
The unramified deformation condition gives a concrete example when
both the residual local representation and \(\delta|_v\) are
unramified: inflate the deformation functor of \(G_{K_v}/I_v\). Its tangent map is the
inflation map of \eqref{v4-arith:coefficients:eq:unramified}.

The tangent coefficients do not usually carry an associative product.
In rank two the trace-zero matrices \(E_{12}\) and \(E_{21}\) have
product \(E_{11}\), whose trace is one. Thus
\(\mathrm{ad}^0\bar\rho\) is not an associative subalgebra of
\(M_2(k)\). A Selmer tangent complex therefore does not acquire the
associative bar construction merely from its deformation-theoretic
name. A separate coefficient algebra, and a comparison respecting its
operations, are needed.

Representability, Noetherianity, perfectness, and cotangent amplitude
are further conditions on the deformation problem. Obstructions in
\(H^2_{\mathcal L}\), which appear in tangent degree one, do not
by themselves rule out quasi-smoothness. For a derived scheme,
quasi-smoothness means that its cotangent complex is perfect of
amplitude \([-1,0]\). This permits tangent degree one.

\begin{example}\label{v4-arith:coefficients:ex:quasismooth}
Let \(A=k[x]/(x^2)\). The element \(x^2\) is a non-zero-divisor
in \(k[x]\), so \(\operatorname{Spec}A\) is a local complete
intersection and is quasi-smooth. Its cotangent complex is
\[
 \mathbb L_{A/k}\simeq
 [A\xrightarrow{\,2x\,}A],
\]
in degrees \(-1,0\). At \(x=0\) the differential vanishes, and
the tangent has \(H^1=k\).

This example works also in characteristic
two. Hence vanishing of the obstruction space is not equivalent to
quasi-smoothness. Stack automorphisms require their own amplitude
convention in addition to this scheme calculation.
\end{example}

\section{Real places and the range of cohomology}
\label{v4-arith:coefficients:sec:real}

Only finite places entered the relative product above. A theory of
compact support over a number field must also specify its complexes
at archimedean places, including any modified Tate-cohomology
convention. At a complex place the absolute Galois group is trivial.
At a real place it is \(C_2\), the group of order two. Coefficients
of characteristic two make the latter distinction unavoidable.

\begin{proposition}\label{v4-arith:coefficients:prop:real-two}
For the trivial \(C_2\)-module \(k=\mathbb F_2\),
\(H^n(C_2,k)=k\) for every \(n\geq0\).
\end{proposition}
\begin{proof}
If \(\sigma\) generates \(C_2\), set \(\eta=\sigma-1\).
Then \(k[C_2]=k[\eta]/(\eta^2)\). A free resolution of the
trivial module is
\[
 \cdots\xrightarrow{\eta}k[C_2]\xrightarrow{\eta}k[C_2]
 \xrightarrow{\eta}k[C_2]\longrightarrow k\longrightarrow0.
\]
The kernel and image of multiplication by \(\eta\) are both
\(k\eta\), so the resolution is exact. Applying
\(\operatorname{Hom}_{k[C_2]}(-,k)\) makes every differential zero.
\end{proof}

Consequently the unmodified real local cochains cannot be assigned a
cohomological bound of two at the coefficient prime two. Neither a
finite-place calculation nor a formal Selmer pullback supplies the
missing archimedean convention. A comparison with de Rham or chiral
objects at infinity must also specify its scalar field and geometric
carrier.

\section{Tame symbols and their recipients}
\label{v4-arith:coefficients:sec:symbols}

Residues provide arithmetic operations with explicitly different
targets at different places. For a field \(F\), the Milnor group
\(K^M_2(F)\) is the abelian group generated by symbols
\(\{a,b\}\), for \(a,b\in F^\times\). Its defining relations
are bilinearity in the multiplicative arguments and the Steinberg
relation \(\{a,1-a\}=0\) for \(a\ne0,1\).

Let \(v:F^\times\to\mathbb Z\)
be a discrete valuation with residue field \(\kappa_v\).
Write \(a=\pi^m u\) and \(b=\pi^n w\), where \(\pi\) is a
uniformizer and \(u,w\) are units. Set
\begin{equation}\label{v4-arith:coefficients:eq:tame}
 \partial_v\{a,b\}=(-1)^{mn}\overline{w}^{m}
                                  \overline{u}^{-n}
 \in\kappa_v^\times.
\end{equation}
This convention gives \(\partial_v\{\pi,u\}=\bar u\).

\begin{proposition}\label{v4-arith:coefficients:prop:tame}
Formula~\eqref{v4-arith:coefficients:eq:tame} is independent of
the uniformizer and defines a homomorphism
\(K^M_2(F)\to\kappa_v^\times\).
\end{proposition}
\begin{proof}
Replacing \(\pi\) by \(\pi t\) replaces the unit parts by
\(ut^{-m}\) and \(wt^{-n}\). The extra powers of \(\bar t\)
cancel. Addition of valuations and multiplication of unit parts
prove bilinearity. It remains to evaluate \(\{a,1-a\}\).

If \(v(a)>0\), then \(1-a\) is a unit of residue one, so the
symbol is one. If \(v(a)=0\), both arguments are units unless
\(v(1-a)>0\). In that case \(\bar a=1\), with the same result.
If \(m=v(a)<0\), then \(v(1-a)=m\) and the ratio of their
unit parts has residue \(-1\). The symbol becomes
\((-1)^{m^2+m}=1\). Thus every Steinberg relation is respected.
\end{proof}

For \(F=\mathbb Q\), the finite symbols land in
\(\bigoplus_p\mathbb F_p^\times\), where all but finitely many
components are one. There is also a real sign symbol
\[
 h_\infty\{a,b\}=
 \begin{cases}-1,&a<0\text{ and }b<0,\\1,&\text{otherwise}.
 \end{cases}
\]
Signs make it bilinear, and \(a\) and \(1-a\) cannot both be
negative, so it too respects the Steinberg relation.

\begin{theorem}[Tate's rational-field calculation]
\label{v4-arith:coefficients:thm:tate}
The combined homomorphism is an isomorphism
\[
 (h_\infty,(\partial_p)_p):K^M_2(\mathbb Q)
 \xrightarrow{\sim}
 \{\pm1\}\oplus\bigoplus_p\mathbb F_p^\times.
\]
\end{theorem}

This is the classical computation of \(K_2(\mathbb Q)\), using
the identification of field \(K_2\) with its Milnor presentation.\footnote{
C. Weibel, \emph{The K-book}, Chapter III, \S6, especially
Theorem 6.5.1 and the real symbol on pp.~231--234.}
The finite symbols alone omit the order-two class \(\{-1,-1\}\):
each finite valuation of its arguments is zero, while its real symbol
is \(-1\). The real factor is therefore necessary even before any
comparison with geometric residues is attempted.

A product formula needs a common recipient. On the rational function
field \(F(t)\), the normed tame symbols satisfy
\begin{equation}\label{v4-arith:coefficients:eq:curve-reciprocity}
 \prod_{x\in|\mathbb P^1_F|}
 N_{\kappa(x)/F}\bigl(\partial_x\xi\bigr)=1
 \quad(\xi\in K^M_2(F(t))).
\end{equation}
Here every factor belongs to \(F^\times\), the point at infinity
is included, and only finitely many factors differ from one.
This is the tame-symbol reciprocity theorem for the projective
line.\footnote{Weibel, \emph{The K-book}, Chapter III,
Example 6.5.3, p.~236.}

For example, if
\(\xi=\{t,u\}\), with \(u\in F^\times\), the symbols at zero
and infinity are \(u\) and \(u^{-1}\). All others are one.

The finite residue fields of a number field do not come with these
norm maps to a single ground field. Multiplying elements of
\(\mathbb F_p^\times\) for varying \(p\) without a specified
recipient is not a reciprocity law. Furthermore,
\(\operatorname{Spec}\mathcal O_K\) has no codimension-two points.
Its dimension-one residue diagram cannot be relabelled as the
two-step flag diagram of an arithmetic surface. The latter requires
an actual surface, its flags, and the corresponding successive residue
and norm maps.

One can choose the incidence poset of finite subsets of finite places,
with inclusions as arrows. It records which places occur together.
It is distinct from the scheme configuration functor of pairwise
disjoint sections of a curve. In particular it supplies neither
collision differential forms nor their operation homotopies.
Comparing the Steinberg and geometric Arnold relations requires a map
between their complexes. Its compatibility with the relevant operations
must also be proved.

\section{Iwasawa coefficients and coalgebra length}
\label{v4-arith:coefficients:sec:iwasawa}

For a procyclic group \(\Gamma\cong\mathbb Z_\ell\), fix a
topological generator \(\gamma\). Its completed group ring is
\[
 \Lambda=\mathbb Z_\ell[[\Gamma]]
 :=\varprojlim_{r,n}R_r[\Gamma/\ell^n\Gamma].
\]
The choice of \(\gamma\) identifies this ring with a power-series
ring. If \(\Gamma\) comes from a cyclotomic extension, this means
its chosen procyclic factor. The full cyclotomic group can also have
a finite factor. In \(\mathbb Z_\ell^\times\), the standard
procyclic subgroup is \(1+\ell\mathbb Z_\ell\) for odd \(\ell\),
and \(1+4\mathbb Z_2\) for \(\ell=2\).

\begin{proposition}\label{v4-arith:coefficients:prop:iwasawa}
The map \(\gamma\mapsto1+T\) induces a topological ring
isomorphism \(\Lambda\cong\mathbb Z_\ell[[T]]\).
\end{proposition}
\begin{proof}
The finite quotients are
\[
 R_r[T]\big/\bigl((1+T)^{\ell^n}-1\bigr).
\]
The ideals \((\ell^r,(1+T)^{\ell^n}-1)\) are cofinal with
the powers of \((\ell,T)\). To see one inclusion, fix a power
\((\ell,T)^N\), take \(r\geq N\), and increase \(n\).
For each \(1\leq j<N\), the identity
\[
 \binom{\ell^n}{j}=\frac{\ell^n}{j}
                         \binom{\ell^n-1}{j-1}
\]
gives valuation at least \(n-v_\ell(j)\). Thus every coefficient
of degree below \(N\) lies in the required power of \(\ell\).

Conversely, modulo the displayed ideal,
\(T^{\ell^n}\) belongs to \(\ell R_r[T]\), because every
intermediate binomial coefficient is divisible by \(\ell\).
Hence \(T^{r\ell^n}=0\), and also \(\ell^r=0\).
A sufficiently large power of \((\ell,T)\) therefore vanishes in
each finite quotient. The two completions agree, and completion of
\(\mathbb Z[T]\) at \((\ell,T)\) is \(\mathbb Z_\ell[[T]]\).
\end{proof}

An augmented dg \(\Lambda\)-algebra has the coefficient bar of
Chapter~\ref{v4-ch:bar-cobar-topology}. Any additional completion must
again name its tensor target. A trace class in Hochschild homology
does not, by itself, provide that algebra or its augmentation.

For a finitely generated torsion \(\Lambda\)-module \(M\), define
\[
 \mu(M)=\operatorname{length}_{\Lambda_{(\ell)}}M_{(\ell)}.
\]
The localization \(\Lambda_{(\ell)}\) is a discrete valuation ring
with uniformizer \(\ell\). This invariant concerns a module.
Conilpotence concerns iterated reduced coproducts of a coaugmented
coalgebra. Even for elementary ordinary \(\Lambda\)-coalgebras,
the two conditions are independent.

\begin{proposition}\label{v4-arith:coefficients:prop:mu}
There is a conilpotent coaugmented \(\Lambda\)-coalgebra whose
coaugmentation quotient has \(\mu=1\). There is also a
nonconilpotent one whose quotient has \(\mu=0\).
\end{proposition}
\begin{proof}
For the first example take
\(C=\Lambda\oplus M\), with \(M=\Lambda/(\ell)\), counit
projection, and every element of \(M\) primitive:
\(\Delta(m)=1\otimes m+m\otimes1\). Its reduced coproduct is
zero, while \(M_{(\ell)}\) has length one.

For the second take \(C=\Lambda\oplus Mx\), with
\(M=\Lambda/(T)\), and set
\[
 \Delta(x)=1\otimes x+x\otimes1+x\otimes x,
 \qquad \epsilon(x)=0.
\]
The relation \(Tx=0\) is respected by this formula.
The element \(1+x\) is group-like, which proves coassociativity.
The counit identities follow directly. Every iterated reduced
coproduct of \(x\) is a nonzero tensor power, because
\(M^{\otimes_\Lambda n}=M\). But \(T\) is invertible in
\(\Lambda_{(\ell)}\), so \(M_{(\ell)}=0\) and \(\mu(M)=0\).
\end{proof}

These are ordinary coalgebras over \(\Lambda\). Their quotient
modules are not required to be flat. They rule out using the module
condition \(\mu=0\) as a definition of conilpotence in this category.
A duality theorem with Iwasawa coefficients must specify its algebra
and coalgebra classes, tensor, and completion. It must then identify
the unit and counit maps it claims to invert.
